% !TeX program = xelatex
% !TeX encoding = UTF-8
% ========================================================================
\documentclass{resume}
% ========================================================================
% ========================================================================
\usepackage{zh_CN_fonts}
\usepackage{linespacing_fix}
\usepackage{amsmath,bm}
\usepackage{booktabs}
\usepackage{cite}
\usepackage[hidelinks,bookmarks=false]{hyperref}
\usepackage{fontawesome} % Icons
\usepackage{datetime2} % Date and time
\usepackage{enumitem} % For custom lists

% ========================================================================
% ========================================================================
\begin{document}
% ========================================================================
% ========================================================================
\pagenumbering{gobble}
% ========================================================================
\begin{figure}[h]
    \begin{minipage}[c]{2.15cm}
        \centering
        \includegraphics[height=2cm,width=2cm]{cydQRcode.jpeg}
    \end{minipage}
    \hfill
    \begin{minipage}[c]{2.5cm}
        \centering
        \includegraphics[height=3.5cm,width=2.5cm]{cyd.jpg}
    \end{minipage}
\end{figure}
\vspace{-28mm}
% ========================================================================
\name{\Large\textbf{Yinda Chen}}
% ========================================================================
\basicInfo{\phone{(+86)13058626611}\quad\email{cyd0806@mail.ustc.edu.cn}
            }

 \centering{\faChain \url{https://ydchen0806.github.io/}} \\
  \centering{\faChain \url{https://scholar.google.com/citations?user=hCvlj5cAAAAJ&hl=en&oi=ao}}
% ========================================================================
\vspace{.5em}

% ========================================================================
% ========================================================================
\section{\makebox[.75em][c]{\faGraduationCap} \textbf{Education}}
% ========================================================================
\datedsubsection{\textbf{University of Science and Technology of China \& Shanghai AI Lab}\ {\faMapMarker \small\textbf{Hefei \& Shanghai}}}{}
\begin{flushleft}
\textbf{Information and Communication Engineering, Ph.D. Candidate} \hfill \textbf{2024.7 \textasciitilde\ 2027.6 (Expected)}
\end{flushleft}
\begin{itemize}
    \item \textbf{Research Focus:}\hangindent=2em
    Predictive multimodal intelligence: self-supervised representation learning, embodied world models and VLA; scientific agents, foundation-model post-training, medical imaging, and connectomics
    \item \textbf{Advisors and Research Mentors:}\hangindent=2em
    \href{https://scholar.google.com/citations?user=5bInRDEAAAAJ&hl=en}{Prof. Feng Wu} (CAE Academician, IEEE Fellow), \href{https://scholar.google.com/citations?user=Snl0HPEAAAAJ&hl=en&oi=ao}{Prof. Zhiwei Xiong}; Co-supervisor (2023): \href{https://scholar.google.com/citations?user=qpBtpGsAAAAJ&hl=en}{Prof. Xiaoou Tang} (IEEE Fellow, Shanghai AI Lab); Connectomics mentor: Ruobing Zhang (Fudan University)
    \item \textbf{Core Courses:}\hangindent=2em
    Algorithm Design and Analysis, Statistical Learning, Deep Learning, Reinforcement Learning
    \item \textbf{Honors:}\hangindent=2em
    NSFC Doctoral Student Project PI (awarded 2024; project period 2025--2027); \href{https://ydchen0806.github.io/docs/cast_ccf_2026_selection_proof.png}{Selected for the 2026 CAST Doctoral Student Program (CCF)}
\end{itemize}

\datedsubsection{\textbf{University of Science and Technology of China}\ {\faMapMarker \small\textbf{Hefei}}}{}
\begin{flushleft}
\textbf{Computer Technology, M.S.} \hfill \textbf{2022.9 \textasciitilde\ 2024.7}
\end{flushleft}
\begin{itemize}
    \item \textbf{Honors:}\hangindent=2em
    National Graduate Scholarship (2022)
\end{itemize}

\datedsubsection{\textbf{Xiamen University}\ {\faMapMarker \small\textbf{Xiamen}}}{}
\begin{flushleft}
\textbf{Environmental Ecological Engineering \& Economics (Dual Degree), B.S.} \hfill \textbf{2018.9 \textasciitilde\ 2022.7}
\end{flushleft}
\begin{itemize}
    \item \textbf{Overall Ranking: 1/31}
    \item \textbf{Honors:}\hangindent=2em
    Xiamen University Academic Star (2021), CDA Level 1 Certification (2022), Kaggle Expert
    \item \textbf{Advisor:}\hangindent=2em
    \href{https://scholar.google.com/citations?user=l1GMXf4AAAAJ&hl=en&oi=ao}{Prof. Yuanye Zhang}
\end{itemize}
% ========================================================================
% ========================================================================
% ========================================================================
% ========================================================================
\section{\makebox[.75em][c]{\faFlask} \textbf{Research Projects}}
\datedsubsection{\textbf{\textit{Embodied Intelligence \& World Models: Pelican Series}}}{\textbf{2026.5 \textasciitilde\ Present}}
\role {Unified Embodied Foundation Models}{Pelican-Unified 1.0 Core Contributor; Pelican-VLA 0.5 Contributor}
\begin{itemize}
    \item In \href{https://arxiv.org/abs/2605.15153}{\textbf{Pelican-Unified 1.0}}, core contributor to \textbf{unified VLM and action modeling}, connecting multimodal understanding, future video prediction, and low-level actions within the unified understanding, reasoning, imagination, and action framework.
    \item Contributed to \href{https://arxiv.org/abs/2607.06655}{\textbf{Pelican-VLA 0.5: Attending Before Acting Benefits Generalization}}, studying \textbf{task-relevant bottleneck representations}, vision-language conditioning, and action prediction, including real-robot closed-loop validation.
    \item Investigated learnable \textbf{Bottleneck Tokens} for task-relevant visual information, the connection between future-state modeling and action learning, and generalization across task scenes.
    \item Media coverage: \href{https://www.peopleapp.com/column/30052155186-500007497131}{People's Daily}, \href{https://www.beijing.gov.cn/fuwu/lqfw/gggs/202605/t20260519_4656998.html}{Beijing Daily / Beijing Municipal Government}, \href{https://mp.weixin.qq.com/s/z-OAqnl4IKHpaiT4bYBHIA}{Machine Heart}, \href{https://www.zhidx.com/p/558258.html}{Zhidongxi}
\end{itemize}

\datedsubsection{\textbf{\textit{NSFC Young Students Basic Research Program (Doctoral Student)}}}{\textbf{2025.1 \textasciitilde \ 2027.12}}
\role {NSFC Project}{Principal Investigator\ {(Awarded 2024; project period 2025--2027; funding \textcolor{red}{300k RMB})}}
\begin{enumerate}[label=\textbf{[\arabic*]}, leftmargin=2em, itemsep=3pt]
    \item \textbf{Neuroscience representations and pretraining}
    \begin{itemize}
        \item Visual token construction, \textbf{autoregressive pretraining}, and multiscale representation consistency for electron microscopy and 3D medical images.
        \item Data-efficient \textbf{self-supervised and predictive learning} with multimodal neural imaging data.
    \end{itemize}
    \item \textbf{Fine-grained neuron segmentation}
    \begin{itemize}
        \item \textbf{Multi-agent masking}, complementary masking consistency, and domain-adaptive segmentation for representation transfer and robustness.
    \end{itemize}
    \item \textbf{Connectome structure--function analysis}
    \begin{itemize}
        \item Investigated neurotransmitter lateralization and its relationship to memory computation in the adult Drosophila mushroom-body connectome.
        \item \textbf{SfN 2026 poster abstract accepted:} Connectome-scale chemical lateralization predicts asymmetric memory computation in the adult Drosophila mushroom body.
    \end{itemize}
\end{enumerate}

\datedsubsection{\textbf{\textit{Representation and Predictive Learning}}}{\textbf{2022.5 \textasciitilde \ Present}}
\role{Conference\&Journal}{Selected first- and co-first-author papers}
\begin{enumerate}[label=\textbf{[\arabic*]}, leftmargin=2em, itemsep=3pt]
    \item \textbf{\href{https://openaccess.thecvf.com/content/ICCV2025/html/Chen_TokenUnify_Scaling_Up_Autoregressive_Pretraining_for_Neuron_Segmentation_ICCV_2025_paper.html}{TokenUnify: Scaling Up Autoregressive Pretraining for Neuron Segmentation}}, \textbf{ICCV 2025} \hfill \mbox{\textcolor{red}{\textbf{[Co-first]}}}
    \begin{itemize}
        \item Developed image \textbf{autoregressive pretraining with Mamba} for transferable neuron-image representations and efficient long-sequence modeling.
        \item Designed \textbf{mixed token-prediction objectives} to capture local structures and long-range dependencies; studied scaling behavior through theoretical analysis and neuron-segmentation experiments.
    \end{itemize}

\begin{samepage}
    \item \textbf{\href{https://www.ijcai.org/proceedings/2023/0068.pdf}{Self-supervised Neuron Segmentation with Multi-agent Reinforcement Learning}}, \textbf{IJCAI 2023} \hfill \mbox{\textcolor{red}{\textbf{[First]}}}
    \begin{itemize}
        \item Developed \textbf{dbMiM}, a multi-agent reinforcement-learning framework that improves masked-image pretraining by jointly optimizing masking ratios and strategies.
        \item Used adaptive masking tasks to learn neuron-image representations for downstream neuron segmentation.
    \end{itemize}
\end{samepage}

    \item \textbf{\href{https://arxiv.org/abs/2308.09917}{Learning Multiscale Consistency for Self-Supervised Electron Microscopy Instance Segmentation}} \\
    \textbf{ICASSP 2024} \hfill \mbox{\textcolor{red}{\textbf{[First]}}}
    \begin{itemize}
        \item Combined \textbf{multiscale feature contrastive learning} and reconstruction for self-supervised EM pretraining.
        \item Designed a \textbf{feature-consistency objective} to address scale-dependent representation differences and evaluated the learned features for neuron instance segmentation.
    \end{itemize}

\end{enumerate}

\datedsubsection{\textbf{\textit{Multimodal and Generative Models}}}{\textbf{2023.8 \textasciitilde \ Present}}
\role{Conference\&Journal}{Selected first- and co-first-author papers}
\begin{enumerate}[label=\textbf{[\arabic*]}, leftmargin=2em, itemsep=3pt]
    \item \textbf{\href{https://doi.org/10.1109/TPAMI.2026.3744812}{Learned Image Coding with Generative Reference of Conditional Latents}} \\
    \textbf{IEEE TPAMI, 2026 (accepted)} \hfill \mbox{\textcolor{red}{\textbf{[Co-first]}}}
    \begin{itemize}
        \item Extended CLC to study \textbf{generative reference images} for conditional latent coding and entropy modeling.
        \item Integrated \textbf{local dictionaries, web retrieval, and image generation} as reference sources, with conditional latent modeling, reference synthesis, and robustness analysis.
    \end{itemize}

    \item \textbf{\href{https://doi.org/10.1609/aaai.v39i12.33403}{Conditional Latent Coding with Learnable Synthesized Reference for Deep Image Compression}} \\
    \textbf{AAAI 2025} \hfill \mbox{\textcolor{red}{\textbf{[Co-first]}}}
    \begin{itemize}
        \item Built a similarity dictionary to retrieve reference images and introduced \textbf{learnable synthesized references} for entropy-model probability estimation.
        \item Combined \textbf{conditional latent coding}, reference synthesis, and adaptive entropy modeling for reference-assisted learned image compression.
    \end{itemize}

    \item \textbf{\href{https://openaccess.thecvf.com/content/ICCV2025W/VLM3D/html/Chen_GTGM_Generative_Text-Guided_3D_Vision-Language_Pretraining_for_Medical_Image_Segmentation_ICCVW_2025_paper.html}{GTGM: Generative Text-Guided 3D Vision-Language Pretraining for Medical Image Segmentation}}, \textbf{ICCV 2025 Workshop (VLM3D)} \hfill \mbox{\textcolor{red}{\textbf{[Co-first]}}}
    \begin{itemize}
        \item Used language-model-generated image descriptions for \textbf{3D vision-language pretraining without real paired text}.
        \item Combined image-text contrastive learning and generated-text guidance, with supervised downstream fine-tuning and evaluation on 10 medical-image segmentation datasets.
    \end{itemize}

    \item \textbf{\href{https://arxiv.org/abs/2403.15992}{BIMCV-R: A Landmark Dataset for 3D CT Text-Image Retrieval}}, \textbf{MICCAI 2024} \hfill \mbox{\textcolor{red}{\textbf{[First]}}}
    \begin{itemize}
        \item Built a 3D CT text-image retrieval dataset containing \textbf{8,069 CT volumes}, supporting retrieval between volumetric images and clinical reports.
        \item Developed the \href{https://github.com/ydchen0806/MedFinder}{\textbf{MedFinder}} alignment and retrieval framework for semantic search of 3D CT data.
    \end{itemize}

    \item \textbf{\href{https://arxiv.org/abs/2412.19080}{MaskFactory: Towards High-quality Synthetic Data Generation for Dichotomous Image Segmentation}} \\
    \textbf{NeurIPS 2024} \hfill \mbox{\textcolor{red}{\textbf{[Co-first]}}}
    \begin{itemize}
        \item Edited masks through \textbf{rigid and non-rigid deformations}, then used ControlNet to generate corresponding mask-image pairs.
        \item Combined deformation-driven mask editing, controlled image generation, and data-quality assessment to study synthetic training data for dichotomous image segmentation.
    \end{itemize}

    \item \textbf{\href{https://ieeexplore.ieee.org/document/11205280}{EMPOWER: Evolutionary Medical Prompt Optimization With Reinforcement Learning}}, \textbf{IEEE JBHI, 2025 (Early Access)} \hfill \mbox{\textcolor{red}{\textbf{[Co-first]}}}
    \begin{itemize}
        \item Combined domain knowledge and \textbf{reinforcement learning} for evolutionary medical prompt optimization, including medical attention, multi-dimensional quality assessment, evolutionary search, and semantic verification.
        \item The reported experiments reduced factual errors by 24.7\% and improved clinician preference by 15.3\% in blinded evaluation.
    \end{itemize}

\end{enumerate}

\begin{samepage}
\datedsubsection{\textbf{\textit{Large Model Engineering}}}{\textbf{2023.9 \textasciitilde \ Present}}
\role {Projects}{Core Member}
\begin{itemize}
    \item \textbf{Foundation-model systems:} worked on image coding, foundation-model training, and LLM-driven quantitative research; used DDP and DeepSpeed for distributed training on \textbf{thousand-GPU clusters}.
    \item \textbf{Embodied and biomedical AI:} contributed to real-robot closed loops, Pelican-Unified pretraining, and embodied training infrastructure on \textbf{100-GPU H200} clusters; conducted medical-image and neuron-model pretraining using \textbf{64-GPU A40} environments.
\end{itemize}
\end{samepage}

% ========================================================================
\section{\makebox[.75em][c]{\faGroup} \textbf{Industry Internship Experience}}
% ========================================================================
\datedsubsection{\textbf{Ubiquant}\hspace{2em} {\faMapMarker
\small\textbf{Beijing}}}{\textbf{2026.5 \textasciitilde \ 2026.10}}
\role {Quantitative Research and Large Language Models}{Research Intern}
\begin{itemize}
    \item Worked on \textbf{quantitative research} and \textbf{general large models}.
    \item Experienced with \textbf{thousand-GPU-scale clusters} and explored LLM applications in factor mining and data-driven decision making.
    \item \textbf{\href{https://arxiv.org/abs/2610.11451}{SAIL}: core contributor; \href{https://arxiv.org/abs/2610.30949}{IQuest-Q1}: co-author} (October 2026 technical reports).
\end{itemize}

\datedsubsection{\textbf{Beijing Innovation Center of Humanoid Robotics (X-Humanoid)}\hspace{2em} {\faMapMarker
\small\textbf{Beijing}}}{\textbf{2025.12 \textasciitilde \ Present}}
\role {Embodied Intelligence World Model Algorithm Team}{Core Member}
\begin{itemize}
    \item Developed \textbf{unified embodied world models} for humanoid robots, connecting task-relevant perception, future-state prediction, and low-level action generation in real-robot closed loops.
    \item Core contributor to unified VLM and action modeling in \textbf{Pelican-Unified 1.0}, including pretraining and planning validation; contributed to task-relevant bottleneck representations and action learning in \textbf{Pelican-VLA 0.5}.
    \item Experienced with \textbf{100-GPU H200} cluster training for large-scale embodied foundation models.
\end{itemize}

\datedsubsection{\textbf{Tencent (IEG)}\hspace{2em} {\faMapMarker
\small\textbf{Shanghai}}}{\textbf{2025.8 \textasciitilde \ 2025.12}}
\role {Qingyun Program}{Research Intern}
\begin{itemize}
    \item Worked on \textbf{game world models} and multimodal scene understanding, including \href{https://main-vla.github.io/}{\textbf{MAIN-VLA}} for Game for Peace, connecting intentions, game states, and actions.
    \item Delivered the \textbf{Honor of Kings Lingbao} project, focusing on \textbf{TTS} and \textbf{real-time match commentary}.
    \item Trained and optimized Lingbao speech generation models using F5-TTS and IndexTTS2.
    \item Researched \textbf{router-free MoE-LoRA}, including linear expert aggregation and merging into the base model.
\end{itemize}

\datedsubsection{\textbf{PLA General Hospital (301 Hospital)}\hspace{2em} {\faMapMarker
\small\textbf{Beijing}}}{\textbf{2023.9 \textasciitilde \ 2024.2}}
\role {Data Compression Group}{Research Intern}
\begin{itemize}
    \item Collaborated with Academician Qionghai Dai's team on efficient data compression research.
    \item Studied CT and MRI compression through content representations, reference images, and efficient coding.
\end{itemize}

% ========================================================================
\section{\makebox[.75em][c]{\faUniversity} \textbf{Academic Internship Experience}}
% ========================================================================
\datedsubsection{\textbf{Imperial College London}\hspace{2em} {\faMapMarker
\small\textbf{London (Remote)}}}{\textbf{2022.11 \textasciitilde \ 2023.8}}
\role {Data Science Institute}{RA}
\begin{itemize}
    \item Worked with Rossella Arcucci and Che Liu on 3D medical multimodal learning, including generated-text-guided pretraining (GTGM) and CT text-image retrieval (BIMCV-R / MedFinder).
    \item Developed image-text contrastive learning and retrieval methods; contributed to \href{https://github.com/ydchen0806/MedFinder}{BIMCV-R}, a dataset containing \textbf{8,069 CT volumes}.
\end{itemize}

\datedsubsection{\textbf{Xiamen University WISER Club}\hspace{2em} {\faMapMarker
\small\textbf{Xiamen}}}{\textbf{2021.8 \textasciitilde \ 2022.7}}
\role {Data Mining Group}{Academic Community Member}
\begin{itemize}
    \item Responsible for designing and discussing data mining courses, teaching clustering and Transformer sections.
    \item Mentored undergraduate machine-learning projects and organized campus competitions.
\end{itemize}

\datedsubsection{\textbf{Wang Yanan Institute for Studies in Economics}\\\textbf{Xiamen University}\hspace{2em} {\faMapMarker
\small\textbf{Xiamen}}}{\textbf{2020.8 \textasciitilde \ 2021.12}}
\role {Econometrics}{Research Assistant}
\begin{itemize}
    \item Assisted Associate Professor Jiong Zhu in completing territorial economic statistics, conducting visual feature extraction of homestead information.
    \item Developed satellite-image analysis tools and extracted visual features to study land-use changes.
\end{itemize}
% ========================================================================
\section{\makebox[.75em][c]{\faTrophy} \textbf{Honors and Awards}}
% ========================================================================
\vspace{.5mm}
\begin{itemize}[leftmargin=0em, itemsep=3pt]
    \item \textbf{\href{https://ydchen0806.github.io/docs/cast_ccf_2026_selection_proof.png}{Selected for the 2026 CAST Doctoral Student Program (CCF)}} \hfill \textbf{2026.10}
    \begin{itemize}[leftmargin=1.5em]
        \item Cultivating society: China Computer Federation (CCF)
    \end{itemize}

    \item \textbf{NSFC Young Students Basic Research Program (Doctoral Student)} \hfill \textbf{2024.12}
    \begin{itemize}[leftmargin=1.5em]
        \item Principal Investigator; awarded 2024, project period 2025--2027, funding RMB 300,000
    \end{itemize}

    \item \textbf{National Graduate Scholarship} \hfill \textbf{2022.12}
    \begin{itemize}[leftmargin=1.5em]
        \item University of Science and Technology of China
    \end{itemize}

    \item \textbf{Xiamen University Academic Star} \hfill \textbf{2021.12}
    \begin{itemize}[leftmargin=1.5em]
        \item Undergraduate research honor
    \end{itemize}

    \item \textbf{'Jingrun Cup' Mathematics Competition (Mathematics Majors)} \hfill \textbf{2021.10}
    \begin{itemize}[leftmargin=1.5em]
        \item First Place, Xiamen University
    \end{itemize}

    \item \textbf{'Internet+' Innovation Competition} \hfill \textbf{2021.08}
    \begin{itemize}[leftmargin=1.5em]
        \item Gold Award, Fujian Province
    \end{itemize}

    \item \textbf{National College Student Mathematics Competition Finals (Non-Mathematics)} \hfill \textbf{2021.05}
    \begin{itemize}[leftmargin=1.5em]
        \item National Second Prize
    \end{itemize}

    \item \textbf{'Challenge Cup' Academic Competition} \hfill \textbf{2021.05}
    \begin{itemize}[leftmargin=1.5em]
        \item First Prize, Fujian Province
    \end{itemize}

    \item \textbf{National College Student Mathematics Competition (Non-Mathematics)} \hfill \textbf{2020.11}
    \begin{itemize}[leftmargin=1.5em]
        \item First Place, Fujian Province
    \end{itemize}
\end{itemize}
% ========================================================================
\section{\makebox[.75em][c]{\faCogs} \textbf{Technical Skills}}
% ========================================================================
\noindent
\begin{minipage}[t]{.5\textwidth}
    \begin{itemize}
    \item \textbf{Programming}: Python, C, C++, Java, MATLAB, \LaTeX
    \item \textbf{Deep Learning}: PyTorch, TensorFlow, DeepSpeed, DDP
    \end{itemize}
\end{minipage}
\begin{minipage}[t]{.5\textwidth}
    \begin{itemize}
    \item \textbf{Languages}: English (TOEFL 110, GRE 328), Chinese (Native)
    \item \textbf{Tools}: Git, Docker, CUDA, HPC
    \end{itemize}
\end{minipage}

% ========================================================================
{\footnotesize\noindent\textbf{Service:} Reviewer for CVPR'25, ICCV'25, WACV'26, NeurIPS'24, ICML'25, ICLR'24, MICCAI'25, ACM MM'24, AISTATS'24, IJCV, and TIP.\par
\noindent\textbf{Invited Talk:} Vision-Language Pretraining Using Generative Methods, JD AI Team, 2023.03.}

% ========================================================================
\end{document}
